\documentclass[10pt,a4paper]{amsart}
\usepackage[margin=19mm]{geometry}
\usepackage{amsmath,amssymb,mathtools}
\usepackage[scr=rsfs,cal=euler]{mathalfa}
\usepackage{xfrac}
\usepackage[unicode,hidelinks,bookmarks=false]{hyperref}
\newcommand{\ie}{i.e.\ }
\newcommand{\cf}{cf.\ }
\newcommand{\resp}{respectively\ }
\newcommand{\Subsection}{\subsection}
\providecommand{\nobreakdash}{\nobreak\hbox{-}}
\setlength{\emergencystretch}{2em}
\multlinegap0pt
\usepackage{mathtools}
\usepackage{amssymb}
\usepackage{xcolor}
\usepackage{dsfont}
\usepackage{tikz}
\usepackage{bm}
\usepackage[shortlabels]{enumitem}
\newtheorem{corollary}{Corollary}
\def\CC{\mathbb{C}}
\def\EE{\mathbb{E}}
\def\cH{\mathcal{H}}
\title{\bfseries The volume of hyperbolic Poisson zero cells:\\ critical divergence and exact second moment}
\author{Tillmann B\"uhler, Christoph Th\"ale}
\date{Source: arXiv:2604.05760v1 (CC BY 4.0)}
\begin{document}
\maketitle

\noindent\emph{Source and licence.} \bfseries The volume of hyperbolic Poisson zero cells:\\ critical divergence and exact second moment. Version arXiv:2604.05760v1 (CC BY 4.0), distributed by arXiv under the Creative Commons Attribution 4.0 International licence, http://creativecommons.org/licenses/by/4.0/, as recorded in the arXiv licence metadata for this exact version and shown on the abstract page https://arxiv.org/abs/2604.05760v1. Full licence notice: \texttt{licenses/arxiv-2604.05760-CC-BY-4.0.txt}.\par\smallskip

\noindent\textit{Editorial scope.} Counted unit: the article's corollary cor:GammaToCritical\_d (source0.tex lines 431--436). The article's complete original proof of it is delivered with it (source0.tex lines 1562--1767, one \texttt{proof} environment of 206 lines, which the article places away from the statement). No mathematics is written, repaired or re-derived: the statement and the proof are the article's own lines, in the article's own notation, and no retained line is substituted or re-flowed.  Retained uncounted as the source-local context the proof needs: lines 1141--1149; lines 1445--1450.  Nothing was omitted from any retained span, so the package carries no omission record.  No retained line contains a citation, so the wrapper carries no bibliography and no bibliography entry.  No retained line contains a figure environment, an image-inclusion command or drawing code, so no image is delivered and the package declares no asset dependency.  Editorial formatting only, in addition: an AMS \texttt{amsart} preamble that restates the article's own declarations and macros used by the retained lines, namely the environments \texttt{corollary} on separate counters, together with the standard packages those lines need.  The retained environments print in this document as Corollary 1 in order of appearance, whereas the article prints the same items with the numbering of its own full document.  The complete native source of the article as distributed in the arXiv e-print for this exact version is bundled unchanged as \texttt{sources/197978/source0.tex}.
\begin{equation}\label{eq:SecondMomentIntermediate_d}
	\EE\cH^d(Z_o)^2
	= \frac{2^{1-d} \pi^{-\frac{d}{2}}}{\Gamma\left(\frac{d}{2}\right)} \left( \frac{\gamma}{4} \pi^{\frac{d-2}{2}} \Gamma\left(\frac{d}{2}\right) \right)^3
	\int_0^\infty
	\frac{\left| \Gamma\left(a + i\frac{\lambda}{2}\right) \right|^6}{
		\left| \Gamma\left(a + \frac{d+1}{2} + i\frac{\lambda}{2}\right) \right|^6}
	\frac{\left| \Gamma\left(\frac{d-1}{2} + i\lambda\right) \right|^2}{\left| \Gamma(i\lambda) \right|^2}
	\, d\lambda.
\end{equation}

    \begin{equation}\label{eq:gamma_ratio}
    \frac{\Gamma(z+\alpha)}{\Gamma(z+\beta)}
	=
	z^{\alpha-\beta}(1+O(1/z)),
	\qquad |z|\to\infty,
    \end{equation}

\begin{corollary}[Critical divergence]\label{cor:GammaToCritical_d}
	As the intensity approaches the critical value from the supercritical side, i.e.\ $\gamma \downarrow \gamma_c^{(d)}$, the second moment of the volume of the zero cell diverges according to
	\[
	\EE\cH^d(Z_o)^2 \sim K_d \left( \gamma - \gamma_c^{(d)} \right)^{-3},\qquad K_d = \frac{\pi^{d+1}}{2^{d-2}} (d-1) \frac{\Gamma\left(\frac{d+1}{2}\right)^2}{\Gamma\left(\frac{d}{2}\right)^4}.
	\]
\end{corollary}

\begin{proof}[Proof of Corollary \ref{cor:GammaToCritical_d}]
	 We start from \eqref{eq:SecondMomentIntermediate_d}, which says that
	\[
	\EE\cH^d(Z_o)^2
	=
	K(d,\gamma)\int_0^\infty R_a(\lambda)\,P_d(\lambda)\,d\lambda,
	\]
	where
	\[
	R_a(\lambda)\coloneqq
	\frac{\left| \Gamma\left(a + i\frac{\lambda}{2}\right) \right|^6}{
		\left| \Gamma\left(a + \frac{d+1}{2} + i\frac{\lambda}{2}\right) \right|^6},
	\qquad
	P_d(\lambda)\coloneqq
	\frac{\left| \Gamma\left(\frac{d-1}{2} + i\lambda\right) \right|^2}{
		\left| \Gamma(i\lambda) \right|^2},
	\]
	and
	\[
	a=\frac{\gamma c_d}{2}+\frac{1-d}{4}
	=\frac{c_d}{2}\bigl(\gamma-\gamma_c^{(d)}\bigr),
	\qquad
	K(d,\gamma)=\frac{\gamma^3\pi^{d-3}}{2^{d+5}}\Gamma\!\left(\frac d2\right)^2.
	\]
	Hence \(\gamma\downarrow \gamma_c^{(d)}\) is equivalent to \(a\downarrow 0\).

	We decompose
	\[
	I(a)\coloneqq\int_0^\infty R_a(\lambda)P_d(\lambda)\,d\lambda
	= I_1(a)+I_2(a),
	\]
	where
	\[
	I_1(a)\coloneqq\int_0^1 R_a(\lambda)P_d(\lambda)\,d\lambda,
	\qquad
	I_2(a)\coloneqq\int_1^\infty R_a(\lambda)P_d(\lambda)\,d\lambda.
	\]

	We first show that the tail term is negligible on the scale \(a^{-3}\).
	By \eqref{eq:gamma_ratio}, uniformly in \(a\in[0,1]\) and \(\lambda\ge1\),
	$
	R_a(\lambda)
	\le C \lambda^{-3(d+1)}
	$.	Moreover, again by \eqref{eq:gamma_ratio},
	\[
	P_d(\lambda)\le C\lambda^{d-1},\qquad \lambda\ge1.
	\]
	Therefore
	\[
	R_a(\lambda)P_d(\lambda)\le C\lambda^{-2d-4},\qquad \lambda\ge1.
	\]
	Since \(\lambda\mapsto \lambda^{-2d-4}\) is integrable on \([1,\infty)\), we obtain
	$\sup_{0<a\le1} I_2(a)<\infty$,
	and hence $a^3 I_2(a)\longrightarrow 0$ as $a\downarrow0$.

	It remains to analyse \(I_1(a)\). Substituting \(\lambda=2au\) yields
	\[
	a^3 I_1(a)=\int_0^{1/(2a)} g_a(u)\,du,
	\qquad
	g_a(u)\coloneqq 2a^4R_a(2au)P_d(2au).
	\]
	Fix \(u\ge0\). Since \(\Gamma(z)\sim 1/z\) as \(z\to0\), we have
	\[
	a^6\bigl|\Gamma(a(1+iu))\bigr|^6\longrightarrow (1+u^2)^{-3}, \qquad a\downarrow 0.
	\]
	Furthermore,
	\[
	\left|\Gamma\!\left(a+\frac{d+1}{2}+iau\right)\right|^6
	\longrightarrow
	\Gamma\!\left(\frac{d+1}{2}\right)^6, \qquad a\downarrow 0,
	\]
	and
	\[
	P_d(2au)\sim 4a^2u^2\,\Gamma\!\left(\frac{d-1}{2}\right)^2,\qquad a\downarrow 0.
	\]
	Consequently,
	\[
	g_a(u)\longrightarrow
	8\,\frac{\Gamma\!\left(\frac{d-1}{2}\right)^2}{
		\Gamma\!\left(\frac{d+1}{2}\right)^6}
	\frac{u^2}{(1+u^2)^3},\qquad a\downarrow 0.
	\]

	To justify dominated convergence, note that on the support of \(g_a\) one has
	\(0\le 2au\le1\). Using \(\Gamma(z)=\Gamma(1+z)/z\), we obtain
	\[
	a\,\Gamma(a(1+iu))
	=
	\frac{\Gamma(1+a(1+iu))}{1+iu},
	\]
	and since \(1+a(1+iu)\) ranges in a compact subset of \(\CC\), it follows that
	\[
	a^6|\Gamma(a(1+iu))|^6\le C'(1+u^2)^{-3}
	\]
    for some constant $C'>0$.
	Moreover,
	\[
	\left|\Gamma\!\left(a+\frac{d+1}{2}+iau\right)\right|
	\]
	is bounded away from \(0\) and \(\infty\), uniformly in \(0<a\le1\) and
	\(0\le u\le 1/(2a)\). Finally, the function
	$
	\lambda\mapsto \frac{P_d(\lambda)}{\lambda^2}
	$
	extends continuously to \(0\), hence is bounded on \([0,1]\), so that
	$
	P_d(2au)\le C a^2u^2
	$ for some constant $C>0$.
	Therefore
	\[
	0\le g_a(u)\mathbf 1_{\{u\le1/(2a)\}}
	\le C\frac{u^2}{(1+u^2)^3},
	\qquad u\ge0,
	\]
	and the right-hand side is integrable on \([0,\infty)\). Thus, by the dominated convergence theorem,
	\[
	a^3 I_1(a)
	\longrightarrow
	8\,\frac{\Gamma\!\left(\frac{d-1}{2}\right)^2}{
		\Gamma\!\left(\frac{d+1}{2}\right)^6}
	\int_0^\infty \frac{u^2}{(1+u^2)^3}\,du.
	\]
	Since
	\[
	\int_0^\infty \frac{u^2}{(1+u^2)^3}\,du=\frac{\pi}{16},
	\]
	we obtain
	\[
	a^3 I(a)\longrightarrow
	\frac{\pi}{2}\,
	\frac{\Gamma\!\left(\frac{d-1}{2}\right)^2}{
		\Gamma\!\left(\frac{d+1}{2}\right)^6},\qquad\text{or equivalently}\qquad I(a)\sim
	\frac{\pi}{2}\,
	\frac{\Gamma\!\left(\frac{d-1}{2}\right)^2}{
		\Gamma\!\left(\frac{d+1}{2}\right)^6}\,
	a^{-3}.
	\]
	Since \(K(d,\gamma)\to K(d,\gamma_c^{(d)})\) as \(\gamma\downarrow\gamma_c^{(d)}\), it follows that
	\[
	\EE\cH^d(Z_o)^2
	\sim
	K(d,\gamma_c^{(d)})
	\frac{\pi}{2}\,
	\frac{\Gamma\!\left(\frac{d-1}{2}\right)^2}{
		\Gamma\!\left(\frac{d+1}{2}\right)^6}\,
	a^{-3}.
	\]
	Using
	\[
	a=\frac{c_d}{2}\bigl(\gamma-\gamma_c^{(d)}\bigr),
	\qquad
	a^{-3}=\frac{8}{c_d^3}\bigl(\gamma-\gamma_c^{(d)}\bigr)^{-3},
	\]
	we arrive at
	\[
	\EE\cH^d(Z_o)^2
	\sim
	\left[
	K(d,\gamma_c^{(d)})
	\frac{4\pi}{c_d^3}
	\frac{\Gamma\!\left(\frac{d-1}{2}\right)^2}{
		\Gamma\!\left(\frac{d+1}{2}\right)^6}
	\right]
	\bigl(\gamma-\gamma_c^{(d)}\bigr)^{-3}.
	\]
	The constant simplifies as follows:
	\[
	K(d,\gamma_c^{(d)}) \frac{4\pi}{c_d^3}
	=
	\frac{\pi^{d-2}}{2^{d+3}} \Gamma\!\left(\frac d2\right)^2
	\left(\frac{\gamma_c^{(d)}}{c_d}\right)^3,
	\]
	and
	\[
	\left(\frac{\gamma_c^{(d)}}{c_d}\right)^3
	=
	8\pi^3(d-1)^3
	\frac{\Gamma\!\left(\frac{d+1}{2}\right)^6}{
		\Gamma\!\left(\frac d2\right)^6}.
	\]
	Hence
	\[
	\EE\cH^d(Z_o)^2
	\sim
	\frac{\pi^{d+1}}{2^d}(d-1)^3
	\frac{\Gamma\!\left(\frac{d-1}{2}\right)^2}{
		\Gamma\!\left(\frac d2\right)^4}
	\bigl(\gamma-\gamma_c^{(d)}\bigr)^{-3}.
	\]
	Finally, using
	$
	\Gamma\!\left(\frac{d-1}{2}\right)
	=
	\frac{2}{d-1}\Gamma\!\left(\frac{d+1}{2}\right),
	$
	we obtain
	\[
	\EE\cH^d(Z_o)^2
	\sim
	\frac{\pi^{d+1}}{2^{d-2}}(d-1)
	\frac{\Gamma\!\left(\frac{d+1}{2}\right)^2}{
		\Gamma\!\left(\frac d2\right)^4}
	\bigl(\gamma-\gamma_c^{(d)}\bigr)^{-3},
	\]
	which proves the claim.
\end{proof}

\end{document}
